\subsection{Parsing}

Using a grammar to generate a string is straightforward.
Parsing solves the opposite problem: is a string part of a
language? Or, is there some combination of rewrites that 
will generate a string, for whatever rewrite are valid 
under the rules?

A parse tree shows how a string was generated.
\begin{itemize}
    \item Interior nodes: non-terminals
    \item Leaf nodes: terminals
    \item Children of interior nodes: the terminals and non-terminals 
    generated by applying a rule
\end{itemize}

Parsing is recognizing members in a language 
specified/defined/generated by a grammar. The parse tree is like 
a sentence diagram. 

In practice, the most common parser structure is top-down,
meaning it functions by expanding the parse tree from the root 
node down. It expands the tree in pre-order, and for each 
node in the parse tree, figures out what it expands to. 
Commonly involves mutually-recursive functions or LL(I) 
parser, a table-based parser that operates similarly 
to recursive-descent. 

